% chapters/ch05.tex
% 第~
\ref{chap:5}~章 \index{Hilbertkongjian@Hilbert 空间 (Hilbert space)}
Hilbert 空间
% 源: 原书 PDF 第 140--165 页 (印刷页 131--156)

\chapter{Hilbert 空间}
\label{chap:5}

\section{定义与初等性质. \index{daotujishangdetouying@到凸集上的投影 (projection on convex set)}
到闭凸集上的投影}
\label{sec:5.1}

\begin{definition}
设 $H$ 为向量空间.\index{biaoliangji@标量积 (scalar product)}
\emph{标量积} $(u,v)$ 是 $H\times H$ 上取值于 $\R$ 的双线性
形式 (即由 $H\times H$ 到 $\R$ 的, 对两个变量都线性的映射), 使得
\begin{align*}
(u,v)&=(v,u) &&\forall u,v\in H\quad(\text{对称性}),\\
(u,u)&\ge0 &&\forall u\in H\quad(\text{正性}),\\
(u,u)&\ne0 &&\forall u\ne0\quad(\text{确定性}).
\end{align*}
\end{definition}

我们回忆, 标量积满足 Cauchy--Schwarz 不等式
\begin{equation*}
\abs{(u,v)}\le(u,u)^{1/2}(v,v)^{1/2}\quad\forall u,v\in H.
\end{equation*}
[记住下面一点有时是有用的: Cauchy--Schwarz 不等式的证明并不需要假设
$(u,u)\ne0\ \forall u\ne0$.] 由 Cauchy--Schwarz 不等式推出, 量
\begin{equation*}
\boxed{\abs{u}=(u,u)^{1/2}}
\end{equation*}
是一个范数——我们将常用 $\abs{\cdot}$ (而不是 $\norm{\cdot}$) 表示由标量积
导出的范数.事实上, 我们有
\begin{equation*}
\abs{u+v}^2=(u+v,u+v)=\abs{u}^2+(u,v)+(v,u)+\abs{v}^2
\le\abs{u}^2+2\abs{u}\abs{v}+\abs{v}^2,
\end{equation*}
从而 $\abs{u+v}\le\abs{u}+\abs{v}$.

我们回忆\index{pingxingsibianxingfaze@平行四边形法则 (parallelogram law)}
经典的\emph{平行四边形法则}:
\begin{equation}
\abs{\frac{a+b}{2}}^2+\abs{\frac{a-b}{2}}^2
=\frac12\bigl(\abs{a}^2+\abs{b}^2\bigr)\quad\forall a,b\in H.
\label{eq:5.1}
\end{equation}

\begin{definition}
\emph{Hilbert 空间}是赋以一个标量积的向量空间 $H$, 使得 $H$ 关于范数
$\abs{\cdot}$ 完备.
\end{definition}

下文 $H$ 总表示一个 Hilbert 空间.

\runin{基本例子.}
赋以标量积
\begin{equation*}
(u,v)=\int_\Omega u(x)v(x)\,d\mu
\end{equation*}
的 $L^2(\Omega)$ 是 Hilbert 空间.特别地, $\ell^2$ 是 Hilbert 空间.第~
\ref{chap:8}~章与
第~
\ref{chap:9}~章中研究的 Sobolev 空间 $H^1$ 是 Hilbert 空间的另一个例子; 它以
$L^2(\Omega)$ 为``模型''.

\begin{keyproposition}
\label{prop:5.1}
$H$ 一致凸, 从而自反.
\end{keyproposition}

\begin{proof}
设 $\varepsilon>0$, $u,v\in H$ 满足 $\abs{u}\le1$, $\abs{v}\le1$ 且
$\abs{u-v}>\varepsilon$.由平行四边形法则有
\begin{equation*}
\abs{\frac{u+v}{2}}^2<1-\frac{\varepsilon^2}{4}
\quad\text{从而}\quad
\abs{\frac{u+v}{2}}<1-\delta,
\quad\text{其中}\ \delta=1-\Bigl(1-\frac{\varepsilon^2}{4}\Bigr)^{1/2}>0.
\end{equation*}
\end{proof}

\begin{keytheorem}[到闭凸集上的投影]
\label{thm:5.2}
设 $K\subset H$ 为非空闭凸集.则对每个 $f\in H$ 存在唯一的元素 $u\in K$,
使得
\begin{equation}
\abs{f-u}=\min_{v\in K}\abs{f-v}=\dist(f,K).
\label{eq:5.2}
\end{equation}
此外, $u$ 由下述性质\emph{刻画}:
\begin{equation}
u\in K\quad\text{且}\quad(f-u,v-u)\le0\quad\forall v\in K.
\label{eq:5.3}
\end{equation}
\end{keytheorem}

\begin{notation}
上述元素 $u$ 称为 $f$ 在 $K$ 上的\emph{投影}, 记作
\begin{equation*}
\boxed{u=P_Kf.}
\end{equation*}
\end{notation}

不等式 \eqref{eq:5.3} 说的是: 向量 $\overrightarrow{uf}$ 与任一向量
$\overrightarrow{uv}$ ($v\in K$) 的标量积 $\le0$, 即这两个向量所成的角
$\theta\ge\pi/2$; 见图  \ref{fig:4}.

\begin{figure}[tbp]
	\centering
	\begin{tikzpicture}[scale=1.0]
		\draw[thick] (0,0) ellipse (2.4 and 1.2);
		\coordinate (u) at (1.62,0.88);
		\coordinate (v) at (-0.35,0.88);
		\coordinate (f) at (2.95,1.78);
		\draw[-{Stealth[length=2mm]}] (u) -- (v);
		\draw[-{Stealth[length=2mm]}] (u) -- (f);
		\fill (u) circle (1.3pt);
		\draw (u) ++(34:0.62) arc[start angle=34, end angle=180, radius=0.62];
		\node at ($(u)+(107:0.95)$) {$\theta$};
		\node at (-0.55,-0.35) {$K$};
		\node[right] at (u) {$u=P_Kf$};
		\node[above left] at (v) {$v$};
		\node[right] at (f) {$f$};
	\end{tikzpicture}
	\caption{}
	\label{fig:4}
\end{figure}

\begin{proof}
\emph{(a) 存在性.} 我们将给出两个不同的证明:

\begin{enumerate}[label=\arabic*., leftmargin=2.2em, itemsep=0.3em]
	\item 函数 $\varphi(v)=\abs{f-v}$ 凸, 连续, 且
	      $\lim_{\abs{v}\to\infty}\varphi(v)=+\infty$.由于 $H$ 自反, 由推论
	       \ref{cor:3.23} 推出 $\varphi$ 在 $K$ 上达到最小值.
	\item 第二个证明\emph{不}依赖自反空间与一致凸空间的理论.它是直接的论证.
	      设 $(v_n)$ 为 \eqref{eq:5.2} 的\emph{极小化序列}, 即 $v_n\in K$ 且
	      \begin{equation*}
	      d_n=\abs{f-v_n}\to d=\inf_{v\in K}\abs{f-v}.
	      \end{equation*}
	      我们断言 $(v_n)$ 是 Cauchy 序列.事实上, 对
	      $a=f-v_n$, $b=f-v_m$ 应用平行四边形法则, 得到
	      \begin{equation*}
	      \abs{f-\frac{v_n+v_m}{2}}^2+\abs{\frac{v_n-v_m}{2}}^2
	      =\frac12\bigl(d_n^2+d_m^2\bigr).
	      \end{equation*}
	      但 $\dfrac{v_n+v_m}{2}\in K$, 从而
	      $\abs{f-\frac{v_n+v_m}{2}}\ge d$.由此推出
	      \begin{equation*}
	      \abs{\frac{v_n-v_m}{2}}^2\le\frac12\bigl(d_n^2+d_m^2\bigr)-d^2
	      \quad\text{且}\quad\lim_{m,n\to\infty}\abs{v_n-v_m}=0.
	      \end{equation*}
	      因此序列 $(v_n)$ 收敛于某个极限 $u\in K$, 且 $d=\abs{f-u}$.
\end{enumerate}

\emph{(b) \eqref{eq:5.2} 与 \eqref{eq:5.3} 的等价性.} 假设 $u\in K$ 满足 \eqref{eq:5.2}, 并设 $w\in K$.我们有
\begin{equation*}
v=(1-t)u+tw\in K\quad\forall t\in[0,1],
\end{equation*}
从而
\begin{equation*}
\abs{f-u}\le\abs{f-[(1-t)u+tw]}=\abs{(f-u)-t(w-u)}.
\end{equation*}
因此
\begin{equation*}
\abs{f-u}^2\le\abs{f-u}^2-2t(f-u,w-u)+t^2\abs{w-u}^2,
\end{equation*}
这蕴含 $2(f-u,w-u)\le t\abs{w-u}^2\ \forall t\in(0,1]$.令 $t\to0$ 即得 \eqref{eq:5.3}.

反之, 假设 $u$ 满足 \eqref{eq:5.3}.则我们有
\begin{equation*}
\abs{u-f}^2-\abs{v-f}^2=2(f-u,v-u)-\abs{u-v}^2\le0\quad\forall v\in K;
\end{equation*}
这蕴含 \eqref{eq:5.2}.

\emph{(c) 唯一性.} 假设 $u_1$ 与 $u_2$ 满足 \eqref{eq:5.3}.我们有
\begin{align}
(f-u_1,v-u_1)&\le0\quad\forall v\in K,\label{eq:5.4}\\
(f-u_2,v-u_2)&\le0\quad\forall v\in K.\label{eq:5.5}
\end{align}
在 \eqref{eq:5.4} 中取 $v=u_2$, 在 \eqref{eq:5.5} 中取 $v=u_1$, 并把相应的两个不等式相加, 得到
$\abs{u_1-u_2}^2\le0$.
\end{proof}

\begin{remark}
\label{rem:5.1}
极小化问题与不等式组相联系, 这并不令人意外.让我们回忆一个熟知的例子.
设 $F:\R\to\R$ 为可微函数, 并设 $u\in[0,1]$ 是 $F$ 在 $[0,1]$ 上达到最小值的
点.则或者 $u\in(0,1)$ 且 $F'(u)=0$, 或者 $u=0$ 且 $F'(u)\le0$, 或者 $u=1$
且 $F'(u)=1$.这三种情形可概括为: $u\in[0,1]$ 且
$F'(u)(v-u)\le0\ \forall v\in[0,1]$; 另见习题 5.10.
% 原书此处印作 "u = 1 and F'(u) = 1", 与极小性不符 (应为 F'(u) >= 0 或
% F'(u) <= 0 的端情形). 两份 OCR 与 PDF 文本层一致, 属原书排印之误.
% 依技能规范: 保留原样并记录, 不擅自更正. 见《项目交接文档》"原书疑似错误".
\end{remark}

\begin{remark}
\label{rem:5.2}
设 $K\subset E$ 为一致凸 Banach 空间 $E$ 中的非空闭凸集.则对每个 $f\in E$
存在唯一的元素 $u\in E$, 使得
\begin{equation*}
\norm{f-u}=\min_{v\in K}\norm{f-v}=\dist(f,K);
\end{equation*}
见习题 3.32.
\end{remark}

\begin{proposition}
\label{prop:5.3}
设 $K\subset H$ 为非空闭凸集.则 $P_K$ 不增大距离, 即
\begin{equation*}
\abs{P_Kf_1-P_Kf_2}\le\abs{f_1-f_2}\quad\forall f_1,f_2\in H.
\end{equation*}
\end{proposition}

\begin{proof}
令 $u_1=P_Kf_1$, $u_2=P_Kf_2$.我们有
\begin{align}
(f_1-u_1,v-u_1)&\le0\quad\forall v\in K,\label{eq:5.6}\\
(f_2-u_2,v-u_2)&\le0\quad\forall v\in K.\label{eq:5.7}
\end{align}
在 \eqref{eq:5.6} 中取 $v=u_2$, 在 \eqref{eq:5.5} 中取 $v=u_1$, 并把相应的两个不等式相加, 得到
\begin{equation*}
\abs{u_1-u_2}^2\le(f_1-f_2,u_1-u_2).
\end{equation*}
由此推出 $\abs{u_1-u_2}\le\abs{f_1-f_2}$.
% 原书此处印作 "v = u1 in (5)", 应为 (7). 依技能规范保留原样并记录.
\end{proof}

\begin{corollary}
\label{cor:5.4}
假设 $M\subset H$ 为闭线性子空间.设 $f\in H$.则 $u=P_Mf$ 由下式刻画:
\begin{equation}
\boxed{u\in M\quad\text{且}\quad(f-u,v)=0\quad\forall v\in M.}
\label{eq:5.8}
\end{equation}
此外, $P_M$ 是线性算子, 称为\emph{正交投影}.
\end{corollary}

\begin{proof}
由 \eqref{eq:5.3} 我们有
\begin{equation*}
(f-u,v-u)\le0\quad\forall v\in M,
\end{equation*}
从而
\begin{equation*}
(f-u,tv-u)\le0\quad\forall v\in M,\ \forall t\in\R.
\end{equation*}
由此推出 \eqref{eq:5.8} 成立.

反之, 若 $u$ 满足 \eqref{eq:5.8}, 则我们有
\begin{equation*}
(f-u,v-u)=0\quad\forall v\in M.
\end{equation*}
$P_M$ 的线性性是显然的.
\end{proof}
\section{Hilbert 空间的对偶空间}
\label{sec:5.2}

在 Hilbert 空间中, 写出连续线性泛函是非常容易的.任取 $f\in H$; 则映射
$u\mapsto(f,u)$ 是 $H$ 上的连续线性泛函.一个值得注意的事实是: $H$ 上的
\emph{一切}连续线性泛函都以这种方式得到:

\index{RieszFrchetbiaoshidingli@Riesz--Fréchet 表示定理 (Riesz--Frechet representation theorem)}
\begin{keytheorem}[Riesz--Fréchet 表示定理]
\label{thm:5.5}
给定任一 $\varphi\in H^\star$, 存在唯一的 $f\in H$, 使得
\begin{equation*}
\dual{\varphi}{u}=(f,u)\quad\forall u\in H.
\end{equation*}
此外,
\begin{equation*}
\abs{f}=\norm{\varphi}_{H^\star}.
\end{equation*}
\end{keytheorem}

\begin{proof}
我们再一次给出两个证明:

\begin{enumerate}[label=\arabic*., leftmargin=2.2em, itemsep=0.3em]
	\item 第一个证明与定理  \ref{thm:4.11} 的证明几乎完全相同.考虑如下定义的
	      映射 $T:H\to H^\star$: 给定任一 $f\in H$, 映射 $u\mapsto(f,u)$ 是
	      $H$ 上的连续线性泛函.它定义了 $H^\star$ 中的一个元素, 记作 $Tf$,
	      使得
	      \begin{equation*}
	      \dual{Tf}{u}=(f,u)\quad\forall u\in H.
	      \end{equation*}
	      显然 $\norm{Tf}_{H^\star}=\abs{f}$.于是 $T$ 是由 $H$ \emph{到}
	      $T(H)$ (它是 $H^\star$ 的闭子空间) 上的线性等距.为完成证明, 只须
	      证明 $T(H)$ 在 $H^\star$ 中稠密.假设 $h$ 是 $H^\star$ 上的连续线性
	      泛函, 且它在 $T(H)$ 上为零.由于 $H$ 自反, $h$ 属于 $H$ 且满足
	      $\dual{Tf}{h}=0\ \forall f\in H$.由此推出 $(f,h)=0\ \forall f\in H$,
	      从而 $h=0$.
	\item 第二个证明是更直接的论证, 它避免了任何对自反性的使用.令
	      $M=\varphi^{-1}(\{0\})$, 于是 $M$ 是 $H$ 的闭子空间.我们总可以假设
	      $M\ne H$ (否则 $\varphi\equiv0$, 定理  \ref{thm:5.5} 的结论是显然的
	      ——只须取 $f=0$).我们断言: 存在某个元素 $g\in H$, 使得
	      \begin{equation*}
	      \abs{g}=1\quad\text{且}\quad(g,v)=0\quad\forall v\in M
	      \quad(\text{从而}\ g\notin M).
	      \end{equation*}
	      事实上, 取 $g_0\in H$ 且 $g_0\notin M$.令 $g_1=P_Mg_0$.则
	      \begin{equation*}
	      g=(g_0-g_1)/\abs{g_0-g_1}
	      \end{equation*}
	      具有所需性质.

	      给定任一 $u\in H$, 令
	      \begin{equation*}
	      v=u-\lambda g\quad\text{其中}\
	      \lambda=\frac{\dual{\varphi}{u}}{\dual{\varphi}{g}}.
	      \end{equation*}
	      注意 $v$ 是有定义的, 因为 $\dual{\varphi}{g}\ne0$; 此外 $v\in M$,
	      因为 $\dual{\varphi}{v}=0$.由此推出 $(g,v)=0$, 即
	      \begin{equation*}
	      \dual{\varphi}{u}=\dual{\varphi}{g}(g,u)\quad\forall u\in H,
	      \end{equation*}
	      取 $f=\dual{\varphi}{g}g$ 即完成证明.
\end{enumerate}
\end{proof}

\begin{keyremark}
\label{rem:5.3}
\emph{$H$ 与 $H^\star$: 等同还是不等同? 三元组 $V\subset H\subset V^\star$.}
定理  \ref{thm:5.5} 断言存在由 $H$ 到 $H^\star$ 上的一个典范等距.因此把
$H$ 与 $H^\star$ 等同起来是``合理的''.我们将\emph{常常}这样做, 但\emph{并非
总是}这样做.下面是一个典型的, 在许多应用中出现的, 应当谨慎对待等同的情形.
设 $H$ 是带有标量积 $(\ ,\ )$ 与相应范数 $\abs{\cdot}$ 的 Hilbert 空间.设
$V\subset H$ 是 $H$ 中的稠密线性子空间.设 $V$ 有自己的范数 $\norm{\cdot}$,
且 $V$ 关于 $\norm{\cdot}$ 是 Banach 空间.设嵌入 $V\subset H$ 连续, 即
\begin{equation*}
\abs{v}\le C\norm{v}\quad\forall v\in V.
\end{equation*}
[例如 $H=L^2(0,1)$, $V=L^p(0,1)$ 且 $p>2$, 或 $V=C([0,1])$.]

存在一个典范映射 $T:H^\star\to V^\star$, 它不过是把 $H$ 上的连续线性泛函
$\varphi$ 限制到 $V$ 上, 即
\begin{equation*}
\dual{T\varphi}{v}_{V^\star,V}=\dual{\varphi}{v}_{H^\star,H}.
\end{equation*}
容易看出 $T$ 具有下述性质:
\begin{enumerate}[label=(\roman*), leftmargin=2.6em, itemsep=0.1em]
	\item $\norm{T\varphi}_{V^\star}\le C\abs{\varphi}_{H^\star}
	      \quad\forall\varphi\in H^\star$,
	\item $T$ 是单射,
	\item 若 $V$ 自反, 则 $R(T)$ 在 $V^\star$ 中稠密.\footnote{然而一般说来
	      $T$ \emph{不是}满射.}
\end{enumerate}

把 $H^\star$ 与 $H$ 等同起来, 并把 $T$ 看作由 $H^\star$ 到 $V^\star$ 的典范
嵌入, 通常写作
\begin{equation}
\boxed{V\subset H\simeq H^\star\subset V^\star,}
\label{eq:5.9}
\end{equation}
其中一切嵌入都连续且稠密 (只要 $V$ 自反).人们说 $H$ 是\emph{枢轴空间}.
注意标量积 $\dual{\cdot}{\cdot}_{V^\star,V}$ 与 $(\ ,\ )$ 只要都有意义就
一致, 即
\begin{equation*}
\dual{f}{v}_{V^\star,V}=(f,v)\quad\forall f\in H,\ \forall v\in V.
\end{equation*}
\end{keyremark}

如果 $V$ 本身是带有自己的标量积 $((\ ,\ ))$ 的 Hilbert 空间, 而该标量积与
范数 $\norm{\cdot}$ 相对应, 则情形变得更微妙.当然, 我们可以借助
$((\ ,\ ))$ 把 $V$ 与 $V^\star$ 等同起来.然而 \eqref{eq:5.9} 就变得荒谬了.这表明
不能同时把 $V$ 与 $H$ 同它们的对偶空间等同起来: 必须作出选择.通常的习惯是
把 $H^\star$ 与 $H$ 等同, 写出 \eqref{eq:5.9}, 而不把 $V^\star$ 与 $V$ 等同 [自然, 仍
存在由 $V$ 到 $V^\star$ 的等距, 但它不被看作恒等映射].下面是一个极有启发的
例子.

令
\begin{equation*}
H=\ell^2=\Bigl\{u=(u_n)_{n\ge1};\ \sum_{n=1}^\infty u_n^2<\infty\Bigr\},
\end{equation*}
赋以标量积 $(u,v)=\sum_{n=1}^\infty u_nv_n$.

令
\begin{equation*}
V=\Bigl\{u=(u_n)_{n\ge1};\ \sum_{n=1}^\infty n^2u_n^2<\infty\Bigr\},
\end{equation*}
赋以标量积 $((u,v))=\sum_{n=1}^\infty n^2u_nv_n$.

显然 $V\subset H$ 且嵌入连续, 并且 $V$ 在 $H$ 中稠密.这里我们把 $H^\star$ 与
$H$ 等同, 而把 $V^\star$ 与空间
\begin{equation*}
V^\star=\Bigl\{f=(f_n)_{n\ge1};\ \sum_{n=1}^\infty\frac{1}{n^2}f_n^2
<\infty\Bigr\}
\end{equation*}
等同, 它比 $H$ 大.标量积 $\dual{\cdot}{\cdot}_{V^\star,V}$ 由
\begin{equation*}
\dual{f}{v}_{V^\star,V}=\sum_{n=1}^\infty f_nv_n
\end{equation*}
给出, 而 Riesz--Fréchet 同构 $T:V\to V^\star$ 由
\begin{equation*}
u=(u_n)_{n\ge1}\mapsto Tu=(n^2u_n)_{n\ge1}
\end{equation*}
给出.

\begin{remark}
\label{rem:5.4}
不借助一致凸空间的理论, 也容易证明 Hilbert 空间自反.只须两次使用
Riesz--Fréchet 同构 (先由 $H$ 到 $H^\star$, 再由 $H^\star$ 到 $H^{\star\star}$).
\end{remark}

\begin{remark}
\label{rem:5.5}
设 $H$ 是与它的对偶空间 $H^\star$ 等同的 Hilbert 空间.设 $M$ 为 $H$ 的
子空间.我们已经在 \S\ref{sec:1.3} 中把 $M^\perp$ 定义为 $H^\star$ 的子空间.现在可以
把它看作 $H$ 的子空间, 即
\begin{equation*}
M^\perp=\{u\in H;\ (u,v)=0\ \forall v\in M\}.
\end{equation*}
显然有 $M\cap M^\perp=\{0\}$.此外, 若 $M$ 闭, 则还有
$M+M^\perp=H$.事实上, 每个 $f\in H$ 都可以写成
\begin{equation*}
f=(P_Mf)+(f-P_Mf)
\end{equation*}
而 $f-P_Mf\in M^\perp$; 更确切地说, $f-P_Mf=P_{M^\perp}f$.

由此推出, 在 Hilbert 空间中每个闭子空间都有补空间 (\S\ref{sec:2.4} 意义下).
\end{remark}
\section{\index{Stampacchiadingli@Stampacchia 定理 (Stampacchia theorem)}
Stampacchia 定理与 \index{LaxMilgramdingli@Lax--Milgram 定理 (Lax--Milgram theorem)}
Lax--Milgram 定理}
\label{sec:5.3}

\begin{definition}
称双线性形式 $a:H\times H\to\R$ 为
\begin{enumerate}[label=(\roman*), leftmargin=2.6em, itemsep=0.25em]
	\item \emph{连续}的, 如果存在常数 $C$, 使得
	      \begin{equation*}
	      \abs{a(u,v)}\le C\abs{u}\abs{v}\quad\forall u,v\in H;
	      \end{equation*}
	\item \emph{强制}的, 如果存在常数 $\alpha>0$, 使得
	      \begin{equation*}
	      a(v,v)\ge\alpha\abs{v}^2\quad\forall v\in H.
	      \end{equation*}
\end{enumerate}
\end{definition}

\begin{theorem}[Stampacchia]
\label{thm:5.6}
假设 $a(u,v)$ 是 $H$ 上的连续强制双线性形式.设 $K\subset H$ 为非空闭凸子集.
则对任意给定的 $\varphi\in H^\star$, 存在唯一的元素 $u\in K$, 使得
\begin{equation}
a(u,v-u)\ge\dual{\varphi}{v-u}\quad\forall v\in K.
\label{eq:5.10}
\end{equation}
此外, 若 $a$ 对称, 则 $u$ 由下述性质刻画:
\begin{equation}
\boxed{u\in K\quad\text{且}\quad
\frac12a(u,u)-\dual{\varphi}{u}
=\min_{v\in K}\Bigl\{\frac12a(v,v)-\dual{\varphi}{v}\Bigr\}.}
\label{eq:5.11}
\end{equation}
\end{theorem}

定理  \ref{thm:5.6} 的证明依赖于下面这个非常经典的结果.

\begin{keytheorem}[Banach 不动点定理——压缩映射原理]
\label{thm:5.7}
设 $X$ 为非空完备度量空间, $S:X\to X$ 为严格压缩, 即
\begin{equation*}
d(Sv_1,Sv_2)\le k\,d(v_1,v_2)\quad\forall v_1,v_2\in X,\ \text{其中}\ k<1.
\end{equation*}
则 $S$ 有唯一不动点, $u=Su$.
\end{keytheorem}

证明见例如 T.~M.~Apostol \bref{Apostol}{1}, G.~Choquet \bref{Choquet}{1}, A.~Friedman \bref{Friedman}{3}.

\begin{proof}[定理 \ref{thm:5.6}]
由 Riesz--Fréchet 表示定理 (定理  \ref{thm:5.5}) 我们知道, 存在唯一的 $f\in H$,
使得
\begin{equation*}
\dual{\varphi}{v}=(f,v)\quad\forall v\in H.
\end{equation*}
另一方面, 若固定 $u\in H$, 则映射 $v\mapsto a(u,v)$ 是 $H$ 上的连续线性泛函.
再一次使用 Riesz--Fréchet 表示定理, 我们找到 $H$ 中某个唯一的元素, 记作
$Au$, 使得 $a(u,v)=(Au,v)\ \forall v\in H$.显然 $A$ 是由 $H$ 到 $H$ 的线性
算子, 满足
\begin{align}
\abs{Au}&\le C\abs{u}\quad\forall u\in H,\label{eq:5.12}\\
(Au,u)&\ge\alpha\abs{u}^2\quad\forall u\in H.\label{eq:5.13}
\end{align}
问题 \eqref{eq:5.10} 归结为寻找某个 $u\in K$, 使得
\begin{equation}
(Au,v-u)\ge(f,v-u)\quad\forall v\in K.
\label{eq:5.14}
\end{equation}
设 $\rho>0$ 为常数 (稍后确定).注意 \eqref{eq:5.14} 等价于
\begin{equation}
(\rho f-\rho Au+u-u,v-u)\le0\quad\forall v\in K,
\label{eq:5.15}
\end{equation}
即
\begin{equation*}
u=P_K(\rho f-\rho Au+u).
\end{equation*}
对每个 $v\in K$, 令 $Sv=P_K(\rho f-\rho Av+v)$.我们断言: 若 $\rho>0$ 选得
适当, 则 $S$ 是严格压缩.事实上, 由于 $P_K$ 不增大距离 (见命题  \ref{prop:5.3}),
我们有
\begin{equation*}
\abs{Sv_1-Sv_2}\le\abs{(v_1-v_2)-\rho(Av_1-Av_2)},
\end{equation*}
从而
\begin{align*}
\abs{Sv_1-Sv_2}^2
&=\abs{v_1-v_2}^2-2\rho(Av_1-Av_2,v_1-v_2)+\rho^2\abs{Av_1-Av_2}^2\\
&\le\abs{v_1-v_2}^2\bigl(1-2\rho\alpha+\rho^2C^2\bigr).
\end{align*}
选取 $\rho>0$ 使得 $k^2=1-2\rho\alpha+\rho^2C^2<1$ (即
$0<\rho<2\alpha/C^2$), 即知 $S$ 有唯一不动点.\footnote{若须数值地计算该不动点,
则取 $\rho=\alpha/C^2$ 是有益的, 这样可使 $k$ 最小, 从而加速 $S$ 的迭代的收敛.}

现在假设形式 $a(u,v)$ 也是对称的.则 $a(u,v)$ 在 $H$ 上定义了一个新的标量积;
相应的范数 $a(u,u)^{1/2}$ 与原范数 $\abs{u}$ 等价.由此推出 $H$ 关于这个新标量积
也是 Hilbert 空间.利用 Riesz--Fréchet 定理, 我们现在可以借助新的标量积来表示
泛函 $\varphi$, 即存在某个唯一的元素 $g\in H$, 使得
\begin{equation*}
\dual{\varphi}{v}=a(g,v)\quad\forall v\in H.
\end{equation*}
问题 \eqref{eq:5.10} 归结为寻找某个 $u\in K$, 使得
\begin{equation}
a(g-u,v-u)\le0\quad\forall v\in K.
\label{eq:5.16}
\end{equation}
\\eqref{eq:5.16} 的解是一位老朋友: $u$ 不过是 $g$ 关于新标量积 $a$ 在 $K$ 上的投影.我们还
知道 (由定理  \ref{thm:5.2}) $u$ 是 $K$ 中使
\begin{equation*}
\min_{v\in K}a(g-v,g-v)^{1/2}
\end{equation*}
达到的唯一元素.这相当于在 $K$ 上极小化函数
\begin{equation*}
v\mapsto a(g-v,g-v)=a(v,v)-2a(g,v)+a(g,g)
=a(v,v)-2\dual{\varphi}{v}+a(g,g),
\end{equation*}
或等价地极小化函数
\begin{equation*}
v\mapsto\frac12a(v,v)-\dual{\varphi}{v}.
\end{equation*}
\end{proof}

\begin{remark}
\label{rem:5.6}
容易验证: 若 $a(u,v)$ 是满足
\begin{equation*}
a(v,v)\ge0\quad\forall v\in H
\end{equation*}
的双线性形式, 则函数 $v\mapsto a(v,v)$ 凸.
\end{remark}

\begin{keycorollary}[Lax--Milgram]
\label{cor:5.8}
假设 $a(u,v)$ 是 $H$ 上的连续强制双线性形式.则对任意给定的
$\varphi\in H^\star$, 存在唯一的元素 $u\in H$, 使得
\begin{equation}
a(u,v)=\dual{\varphi}{v}\quad\forall v\in H.
\label{eq:5.17}
\end{equation}
此外, 若 $a$ 对称, 则 $u$ 由下述性质刻画:
\begin{equation}
\boxed{u\in H\quad\text{且}\quad
\frac12a(u,u)-\dual{\varphi}{u}
=\min_{v\in H}\Bigl\{\frac12a(v,v)-\dual{\varphi}{v}\Bigr\}.}
\label{eq:5.18}
\end{equation}
\end{keycorollary}

\begin{proof}
在定理  \ref{thm:5.6} 中取 $K=H$, 并仿照推论  \ref{cor:5.4} 的证明进行论证.
\end{proof}

\begin{remark}
\label{rem:5.7}
Lax--Milgram 定理是求解线性椭圆型偏微分方程的一个非常简单而有效的工具 (见
第~
\ref{chap:8}~章与第~
\ref{chap:9}~章).注意方程 \eqref{eq:5.17} 与极小化问题 \eqref{eq:5.18} 之间的联系是很有意思的.
当这类问题在力学或物理学中出现时, 它们往往有自然的解释: 最小作用量原理, 能量的极小化, 等等.用变分法的语言来说, \eqref{eq:5.17} 是与极小化问题 \eqref{eq:5.18} 相联系的
Euler 方程.粗略地说, \eqref{eq:5.17} 说的是``$F'(u)=0$'', 其中 $F$ 是函数
$F(v)=\frac12a(v,v)-\dual{\varphi}{v}$.
\end{remark}

\begin{remark}
\label{rem:5.8}
有一个直接而初等的论证可以证明 \eqref{eq:5.17} 有唯一解.事实上, 这相当于证明
\begin{equation*}
\forall f\in H\ \ \exists u\in H\ \text{唯一, 使得}\ Au=f,
\end{equation*}
即 $A$ 是由 $H$ 到 $H$ 上的双射.这是下述事实的平凡推论:
\begin{enumerate}[label=(\alph*), leftmargin=2.2em, itemsep=0.2em]
	\item $A$ 是单射 (因为 $A$ 强制),
	\item $R(A)$ 闭, 因为 $\alpha\abs{v}\le\abs{Av}\ \forall v\in H$
	      (强制性的推论),
	\item $R(A)$ 稠密; 事实上, 假设 $v\in H$ 满足
	      \begin{equation*}
	      (Au,v)=0\quad\forall u\in H,
	      \end{equation*}
	      则 $v=0$.
\end{enumerate}
\end{remark}

\section{Hilbert 和. 正交基}
\label{sec:5.4}

\begin{definition}
设 $(E_n)_{n\ge1}$ 为 $H$ 的一列闭子空间.若
\begin{enumerate}[label=(\alph*), leftmargin=2.2em, itemsep=0.25em]
	\item 空间 $E_n$ 两两正交, 即
	      \begin{equation*}
	      (u,v)=0\quad\forall u\in E_n,\ \forall v\in E_m,\ m\ne n,
	      \end{equation*}
	\item 由 $\bigcup_{n=1}^\infty E_n$ 张成的线性空间在 $H$ 中稠密,\footnote{由
	      $E_n$ 张成的线性空间按代数意义理解, 即属于诸空间 $(E_n)$ 的元素的有
	      限线性组合.}
\end{enumerate}
则说 $H$ 是诸 $E_n$ 的\emph{Hilbert 和}, 并记作 $H=\bigoplus_nE_n$.
\end{definition}

\begin{keytheorem}
\label{thm:5.9}
假设 $H$ 是诸 $E_n$ 的 Hilbert 和.给定 $u\in H$, 令
\begin{equation*}
u_n=P_{E_n}u\quad\text{及}\quad S_n=\sum_{k=1}^n u_k.
\end{equation*}
则我们有
\begin{equation}
\lim_{n\to\infty}S_n=u
\label{eq:5.19}
\end{equation}
以及
\begin{equation}
\sum_{k=1}^\infty\abs{u_k}^2=\abs{u}^2
\quad(\text{Bessel--Parseval 恒等式}).
\label{eq:5.20}
\end{equation}
\end{keytheorem}

使用下述引理是方便的.

\begin{lemma}
\label{lem:5.1}
假设 $(v_n)$ 为 $H$ 中任一序列, 使得
\begin{align}
(v_m,v_n)&=0\quad\forall m\ne n,\label{eq:5.21}\\
\sum_{k=1}^\infty\abs{v_k}^2&<\infty.\label{eq:5.22}
\end{align}
令
\begin{equation*}
S_n=\sum_{k=1}^n v_k.
\end{equation*}
则
\begin{equation*}
S=\lim_{n\to\infty}S_n
\end{equation*}
存在, 并且
\begin{equation}
\abs{S}^2=\sum_{k=1}^\infty\abs{v_k}^2.
\label{eq:5.23}
\end{equation}
\end{lemma}

\begin{proof}[引理 \ref{lem:5.1}]
注意对 $m>n$ 有
\begin{equation*}
\abs{S_m-S_n}^2=\sum_{k=n+1}^m\abs{v_k}^2.
\end{equation*}
由此推出 $S_n$ 是 Cauchy 序列, 从而 $S=\lim_{n\to\infty}S_n$ 存在.另一方面,
我们有
\begin{equation*}
\abs{S_n}^2=\sum_{k=1}^n\abs{v_k}^2.
\end{equation*}
令 $n\to\infty$ 即得 \eqref{eq:5.23}.
\end{proof}

\begin{proof}[定理 \ref{thm:5.9}]
由于 $u_n=P_{E_n}u$, 我们有 (由 \eqref{eq:5.8})
\begin{equation}
(u-u_n,v)=0\quad\forall v\in E_n,
\label{eq:5.24}
\end{equation}
特别地,
\begin{equation*}
(u,u_n)=\abs{u_n}^2.
\end{equation*}
把这些等式相加, 得到
\begin{equation*}
(u,S_n)=\sum_{k=1}^n\abs{u_k}^2.
\end{equation*}
但我们还有
\begin{equation}
\sum_{k=1}^n\abs{u_k}^2=\abs{S_n}^2,
\label{eq:5.25}
\end{equation}
从而得到
\begin{equation*}
(u,S_n)=\abs{S_n}^2.
\end{equation*}
由此推出 $\abs{S_n}\le\abs{u}$, 因此 $\sum_{k=1}^n\abs{u_k}^2\le\abs{u}^2$.

于是我们可以应用引理  \ref{lem:5.1}, 断言 $S=\lim_{n\to\infty}S_n$ 存在.让我们
\emph{不}假设 (b) 而确定 $S$.设 $F$ 为由诸 $E_n$ 张成的线性空间.我们断言
\begin{equation}
S=P_{\overline{F}}u.
\label{eq:5.26}
\end{equation}
事实上, 我们有
\begin{equation*}
(u-S_n,v)=0\quad\forall v\in E_m,\ m\le n
\end{equation*}
(只须写出 $u-S_n=(u-u_m)-\sum_{k\ne m}u_k$).令 $n\to\infty$ 得到
\begin{equation*}
(u-S,v)=0\quad\forall v\in E_m,\ \forall m,
\end{equation*}
从而
\begin{equation*}
(u-S,v)=0\quad\forall v\in F,
\end{equation*}
这蕴含
\begin{equation*}
(u-S,v)=0\quad\forall v\in\overline{F}.
\end{equation*}
另一方面, $S_n\in F\ \forall n$, 而取极限得 $S\in\overline{F}$.这就证明了 \eqref{eq:5.26}.
当然, 若 (b) 成立, 则 $\overline{F}=H$, 从而 $S=u$.在 \eqref{eq:5.25} 中令
$n\to\infty$ 取极限, 即得 \eqref{eq:5.20}.
\end{proof}

\begin{definition}
称 $H$ 中的序列 $(e_n)_{n\ge1}$ 为 $H$ 的\emph{正交基} (或\emph{Hilbert
基},\footnote{不要与代数 (= Hamel) 基混淆; 后者是 $H$ 中的族 $(e_i)_{i\in I}$,
使得每个 $u\in H$ 都能唯一地写成诸 $e_i$ 的有限线性组合 (见习题 1.5).}或
在没有混淆时简称为\emph{基}\footnote{有些作者说 $(e_n)$ 是完全正交系.}), 如果
它满足下述性质:
\begin{enumerate}[label=(\roman*), leftmargin=2.6em, itemsep=0.2em]
	\item $\abs{e_n}=1\ \forall n$ 且 $(e_m,e_n)=0\ \forall m\ne n$,
	\item 由诸 $e_n$ 张成的线性空间在 $H$ 中稠密.
\end{enumerate}
\end{definition}

\begin{keycorollary}
\label{cor:5.10}
设 $(e_n)$ 为正交基.则对每个 $u\in H$, 我们有
\begin{equation*}
u=\sum_{k=1}^\infty(u,e_k)e_k,
\quad\text{即}\quad u=\lim_{n\to\infty}\sum_{k=1}^n(u,e_k)e_k
\end{equation*}
以及
\begin{equation*}
\abs{u}^2=\sum_{k=1}^\infty\abs{(u,e_k)}^2.
\end{equation*}
反之, 给定任一序列 $(\alpha_n)\in\ell^2$, 级数 $\sum_{k=1}^\infty\alpha_ke_k$
收敛到某个元素 $u\in H$, 使得 $(u,e_k)=\alpha_k\ \forall k$ 且
$\abs{u}^2=\sum_{k=1}^\infty\alpha_k^2$.
\end{keycorollary}

\begin{proof}
注意 $H$ 是诸空间 $E_n=\R e_n$ 的 Hilbert 和, 且 $P_{E_n}u=(u,e_n)e_n$.
使用定理  \ref{thm:5.9} 与引理  \ref{lem:5.1}.
\end{proof}

\begin{remark}
\label{rem:5.9}
一般说来, 定理  \ref{thm:5.9} 中的级数 $\sum u_k$ 与推论  \ref{cor:5.10} 中的
级数 $\sum(u,e_k)e_k$ \emph{不是}绝对收敛的, 即可能出现
$\sum_{k=1}^\infty\abs{u_k}=\infty$ 或
$\sum_{k=1}^\infty\abs{(u,e_k)}=\infty$.
\end{remark}

\begin{keytheorem}
\label{thm:5.11}
每个可分 Hilbert 空间都有正交基.
\end{keytheorem}

\begin{proof}
设 $(v_n)$ 为 $H$ 的可数稠密子集.用 $F_k$ 表示由
$\{v_1,v_2,\dots,v_k\}$ 张成的线性空间.序列 $(F_k)$ 是有限维空间的不减序列,
使得 $\bigcup_{k=1}^\infty F_k$ 在 $H$ 中稠密.在 $F_1$ 中任取一个单位向量
$e_1$.若 $F_2\ne F_1$, 则在 $F_2$ 中存在某个向量 $e_2$, 使得 $\{e_1,e_2\}$ 是
$F_2$ 的正交基.重复同样的构造, 即得 $H$ 的一个正交基.
\end{proof}

\begin{remark}
\label{rem:5.10}
定理  \ref{thm:5.11} 与推论  \ref{cor:5.10} 合起来表明: 一切可分 Hilbert 空间都
与空间 $\ell^2$ 同构且等距.尽管这个结果看似惊人, 考虑诸如 $L^2(\Omega)$
(或 Sobolev 空间 $H^1(\Omega)$ 等) 这样的其他 Hilbert 空间仍然是十分重要的.
原因是: 许多很好的线性 (或非线性) 算子, 一旦用基写出来, 可能显得很可怕.
\end{remark}

\begin{remark}
\label{rem:5.11}
若 $H$ 是不可分的 Hilbert 空间——这是一种相当不寻常的情形——则仍可 (借助
Zorn 引理) 证明不可数正交基 $(e_i)_{i\in I}$ 的存在性; 见例如 W.~Rudin
\bref{Rudin}{2}, A.~E.~Taylor--D.~C.~Lay \bref{TaylorLay}{1}, G.~B.~Folland \bref{Folland}{2}, G.~Choquet \bref{Choquet}{1}.
\end{remark}
\section*{第 5 章评注}
\addcontentsline{toc}{section}{第 5 章评注}
\markright{第 5 章评注}

\cmtitem{1. Hilbert 空间的刻画.}
有时知道一个给定的范数 $\norm{\cdot}$ 在向量空间 $E$ 上是否为 Hilbert 范数
是有用的, 即是否存在 $E$ 上的标量积 $(\ ,\ )$, 使得
$\norm{u}=(u,u)^{1/2}\ \forall u\in E$.已知若干种判别准则:

\begin{cmtlettertheorem}{a}{Fréchet--von Neumann--Jordan}
\label{thm:5.12}
假设范数 $\norm{\cdot}$ 满足平行四边形法则 \eqref{eq:5.1}.则 $\norm{\cdot}$ 是 Hilbert
范数.
\end{cmtlettertheorem}

证明见 K.~Yosida \bref{Yosida}{1} 或习题 5.1.

\begin{cmtlettertheorem}{b}{Kakutani}
\label{thm:5.13}
假设 $E$ 是 $\dim E\ge3$ 的赋范空间.假设每个二维子空间 $F$ 都有一个范数为
1 的投影算子 (即存在有界线性投影算子 $P:E\to F$, 使得 $Pu=u\ \forall u\in F$
且 $\norm{P}\le1$).\footnote{我们指出, 每个一维子空间总有范数为 1 的投影算子.
(使用 Hahn--Banach 定理.)}则 $\norm{\cdot}$ 是 Hilbert 范数.
\end{cmtlettertheorem}

\begin{cmtlettertheorem}{c}{de Figueiredo--Karlovitz}
\label{thm:5.14}
设 $E$ 是 $\dim E\ge3$ 的赋范空间.考虑单位球上的径向投影, 即
\begin{equation*}
Tu=\begin{cases}
u & \text{若}\ \norm{u}\le1,\\
u/\norm{u} & \text{若}\ \norm{u}>1.
\end{cases}
\end{equation*}
假设\footnote{可以证明, 在任意赋范空间中, $T$ 满足
$\norm{Tu-Tv}\le2\norm{u-v}\ \forall u,v\in E$, 且常数 2 不能改进;
见习题 5.6.}
\begin{equation*}
\norm{Tu-Tv}\le\norm{u-v}\quad\forall u,v\in E.
\end{equation*}
则 $\norm{\cdot}$ 是 Hilbert 范数.
\end{cmtlettertheorem}

最后, 让我们回忆一个已经提到过的结果 (注记 2.8).

\begin{cmtlettertheorem}{d}{Lindenstrauss--Tzafriri}
\label{thm:5.15}
假设 $E$ 是 Banach 空间, 使得每个闭子空间都有补空间.\footnote{这等价于说
每个闭子空间都有一个有界投影算子 $P$.注意这里——与定理  \ref{thm:5.13}
相反——我们不假设 $\norm{P}\le1$.}则 $E$ 可 Hilbert 化, 即存在一个等价的
Hilbert 范数.
\end{cmtlettertheorem}

\cmtitem{2. 变分不等式.}
Stampacchia 定理是变分不等式理论的出发点 (见例如
D.~Kinderlehrer--G.~Stampacchia \bref{KinderlehrerStampacchia}{1}), 该理论在
力学与物理学 (见例如 G.~Duvaut--J.~L.~Lions \bref{DuvautLions}{1}), 自由边界
值问题 (见例如 C.~Baiocchi--A.~Capelo \bref{BaiocchiCapelo}{1} 与
A.~Friedman \bref{Friedman}{4}), 最优控制 (见例如 J.-L.~Lions
\bref{Lions}{2} 与 V.~Barbu \bref{Barbu}{2}), 随机控制 (见
A.~Bensoussan--J.-L.~Lions \bref{BensoussanLions}{1}) 中有大量应用.

\cmtitem{3. 与单调算子相联系的非线性方程.}
Stampacchia 定理与 Lax--Milgram 定理可以推广到某些类非线性算子.例如,
让我们提到下述结果.

\begin{theorem}[Minty--Browder]
\label{thm:5.16}
设 $E$ 为自反 Banach 空间.设 $A:E\to E^\star$ 为连续非线性映射, 使得
\begin{equation*}
\dual{Av_1-Av_2}{v_1-v_2}>0\quad\forall v_1,v_2\in E,\ v_1\ne v_2,
\end{equation*}
且
\begin{equation*}
\lim_{\norm{v}\to\infty}\frac{\dual{Av}{v}}{\norm{v}}=\infty.
\end{equation*}
则对每个 $f\in E^\star$ 存在方程 $Au=f$ 的唯一解 $u\in E$.
\end{theorem}

有兴趣的读者可在 F.~Browder \bref{Browder}{1} 与 J.-L.~Lions
\bref{Lions}{3} 中找到定理  \ref{thm:5.16} 的证明以及许多推广与应用; 另见
问题 31.

\cmtitem{4. 特殊的正交基. Fourier 级数. 小波.}
在第~
\ref{chap:6}~章中我们将介绍一种构造正交基的非常有力的方法, 即取紧自伴算子的
特征向量.在实践中, 人们常常使用 $L^2(\Omega)$ 的由微分算子的特征函数组成的
特殊基 (见 \S\ref{sec:8.6} 与 \S\ref{sec:9.8}).由
\begin{equation*}
e_n(x)=\sqrt{2/\pi}\sin nx,\quad n\ge1,
\end{equation*}
或
\begin{equation*}
e_n(x)=\sqrt{2/\pi}\cos nx,\quad n\ge0,
\end{equation*}
定义的 $L^2(0,\pi)$ 上的正交基很受喜爱, 因为它导致 Fourier 级数与调和分析
——一个本身就很重要的领域; 见例如 J.~M.~Ash \bref{Ash}{1}, H.~Dym--H.~P.~McKean \bref{DymMcKean}{1}, Y.~Katznelson
\bref{Katznelson}{1}, C.~S.~Rees--S.~M.~Shah--C.~V.~Stanojevic
\bref{ReesShahStanojevic}{1}.

这里有一个使分析学家困惑了数十年的问题.给定 $u\in L^2(0,\pi)$, 考虑它的
Fourier 级数 $S_n=\sum_{k=1}^n(u,e_k)e_k$.已知 (见推论  \ref{cor:5.10})
$S_n\to u$ 在 $L^2(0,\pi)$ 中.由此推出子列 $S_{n_k}\to u$ 在 $(0,\pi)$ 上
几乎处处 (见定理  \ref{thm:4.9}).但是否可以说整个序列 $S_n\to u$ 在
$(0,\pi)$ 上几乎处处? 答案由下述非常深刻的结果给出:

\begin{theorem}[Carleson]
\label{thm:5.17}
若 $u\in L^2(0,\pi)$, 则 $S_n\to u$ 几乎处处.
\end{theorem}

$L^2(0,1)$ 或 $L^2(\R)$ 的其他经典基与 Bessel, Legendre, Hermite, Laguerre, Chebyshev, Jacobi 等人的名字相联系.我们把有兴趣的读者引向
R.~Courant--D.~Hilbert \bref{CourantHilbert}{1} 第 1 卷以及
R.~Dautray--J.-L.~Lions \bref{DautrayLions}{1} 第 VIII 章; 另见第~
\ref{chap:8}~章末尾的
评注 (Sturm--Liouville 算子的谱性质).近来, 由阶梯函数组成的 $L^2(0,1)$ 的
Haar 基与 Walsh 基也引起了很大兴趣; 见例如习题 5.31, 5.32, G.~Alexits
\bref{Alexits}{1}, H.~F.~Harmuth \bref{Harmuth}{1}.

小波理论提供了一种非常重要而优美的新型基.它是分解函数, 信号, 语音, 图像等的
有力工具.有兴趣的读者可以查阅 Y.~Meyer \bref{Meyer}{1}, \bref{Meyer}{2}, \bref{Meyer}{3}, R.~Coifman 与 Y.~Meyer \bref{CoifmanMeyer}{1}, I.~Daubechies \bref{Daubechies}{1}, G.~David \bref{David}{1}, C.~K.~Chui \bref{Chui}{1}, M.~B.~Ruskai 等 \bref{RuskaiEtAl}{1}, J.~J.~Benedetto--M.~W.~Frazier \bref{BenedettoFrazier}{1}, G.~Kaiser
\bref{Kaiser}{1}, J.~P.~Kahane--P.~G.~Lemarié-Rieusset
\bref{KahaneLemarie}{1}, S.~Mallat \bref{Mallat}{1}, G.~Bachman--L.~Narici--E.~Beckenstein \bref{BachmanNariciBeckenstein}{1}, T.~F.~Chan--J.~Shen \bref{ChanShen}{1}, P.~Wojtaszczyk \bref{Wojtaszczyk}{1}, E.~Hernandez--G.~Weiss \bref{HernandezWeiss}{1} 等人的近著及其中的参考文献.

\cmtitem{5. Banach 空间中的 Schauder 基.}
设 $E$ 为 Banach 空间.称序列 $(e_n)_{n\ge1}$ 为\emph{Schauder 基}, 如果对每个
$u\in E$ 存在唯一的序列 $(\alpha_n)_{n\ge1}\subset\R$, 使得
$u=\sum_{k=1}^\infty\alpha_ke_k$ (即
$u=\lim_{n\to\infty}\sum_{k=1}^n\alpha_ke_k$).这样的基在 Banach 空间的几何
理论中起着重要作用 (见例如 B.~Beauzamy \bref{Beauzamy}{1}, J.~Lindenstrauss--L.~Tzafriri \bref{LindenstraussTzafriri}{2}, J.~Diestel \bref{Diestel}{2}, R.~C.~James \bref{James}{2}).分析中使用的
一切经典 (可分) Banach 空间都有 Schauder 基 (见例如 I.~Singer
\bref{Singer}{1}).这一事实使 Banach 猜测每个可分 Banach 空间都有基.经过
数十年徒劳的努力之后, P.~Enflo \bref{Enflo}{1} 发现了反例.甚至可以在
$\ell^p$ ($1<p<\infty$, $p\ne2$) 中构造没有 Schauder 基的闭子空间 (见
J.~Lindenstrauss--L.~Tzafriri \bref{LindenstraussTzafriri}{2}).A.~Szankowski
\bref{Szankowski}{1} 发现了另一个惊人的例子: 当 $H$ 为无限维可分 Hilbert
空间时, $\mathcal{L}(H)$ (配以通常的范数) 没有 Schauder 基.在第~
\ref{chap:6}~章中我们
将看到, 关于紧算子的一个相关问题也有否定的答案.
\section*{第 5 章习题}
\addcontentsline{toc}{section}{第 5 章习题}
\markright{第 5 章习题}

下文 $H$ 总表示赋以标量积 $(\ ,\ )$ 及相应范数 $\abs{\cdot}$ 的 Hilbert 空间.

\begin{exercise}[平行四边形法则.]
\label{exer:5.1}
假设 $E$ 是赋以范数 $\norm{\cdot}$ 的向量空间, 该范数满足平行四边形法则, 即
\begin{equation*}
\norm{a+b}^2+\norm{a-b}^2=2\bigl(\norm{a}^2+\norm{b}^2\bigr)
\quad\forall a,b\in E.
\end{equation*}
我们的目的是证明由
\begin{equation*}
(u,v)=\frac12\bigl(\norm{u+v}^2-\norm{u}^2-\norm{v}^2\bigr),
\quad u,v\in E,
\end{equation*}
定义的量是标量积, 且 $(u,u)=\norm{u}^2$.
\begin{enumerate}[label=\arabic*., leftmargin=2.2em, itemsep=0.3em]
	\item 验证
	      \begin{equation*}
	      (u,v)=(v,u),\quad(-u,v)=-(u,v)\quad\text{且}\quad(u,2v)=2(u,v)
	      \quad\forall u,v\in E.
	      \end{equation*}
	\item 证明
	      \begin{equation*}
	      (u+v,w)=(u,w)+(v,w)\quad\forall u,v,w\in E.
	      \end{equation*}
	      [提示: 依次对 (i) $a=u$, $b=v$; (ii) $a=u+w$, $b=v+w$;
	      (iii) $a=u+v+w$, $b=w$ 使用平行四边形法则.]
	\item 证明 $(\lambda u,v)=\lambda(u,v)\ \forall\lambda\in\R,\ \forall u,v\in E$.

	      [提示: 先考虑 $\lambda\in\N$ 的情形, 再考虑 $\lambda\in\Q$, 最后
	      考虑 $\lambda\in\R$.]
	\item 作出结论.
\end{enumerate}
\end{exercise}

\begin{exercise}[当 $p\ne2$ 时 $L^p$ 不是 Hilbert 空间.]
\label{exer:5.2}
设 $\Omega$ 为测度空间, 并假设存在可测集 $A\subset\Omega$, 使得
$0<\abs{A}<\abs{\Omega}$.

证明: 对任意 $1\le p\le\infty$, $p\ne2$, 范数 $\norm{\cdot}_p$ 都不满足
平行四边形法则.

[提示: 使用支集互不相交的函数.]
\end{exercise}

\begin{exercise}
\label{exer:5.3}
设 $(u_n)$ 为 $H$ 中的序列, $(t_n)$ 为 $(0,\infty)$ 中的序列, 使得
\begin{equation*}
(t_nu_n-t_mu_m,u_n-u_m)\le0\quad\forall m,n.
\end{equation*}
\begin{enumerate}[label=\arabic*., leftmargin=2.2em, itemsep=0.3em]
	\item 假设序列 $(t_n)$ 不减 (可能无界).证明序列 $(u_n)$ 收敛.

	      [提示: 证明序列 $(\abs{u_n})$ 不增.]
	\item 假设序列 $(t_n)$ 不增.证明下述二者必居其一:
	      \begin{enumerate}[label=(\roman*), leftmargin=2.4em, itemsep=0.1em]
	      	\item 或者 $\abs{u_n}\to\infty$,
	      	\item 或者 $(u_n)$ 收敛.
	      \end{enumerate}
	      若 $t_n\to t>0$, 证明 $(u_n)$ 收敛; 而若 $t_n\to0$, 证明情形 (i)
	      与 (ii) 都可能出现.
\end{enumerate}
\end{exercise}

\begin{exercise}
\label{exer:5.4}
设 $K\subset H$ 为非空闭凸集.设 $f\in H$, $u=P_Kf$.证明
\begin{equation*}
\abs{v-u}^2\le\abs{v-f}^2-\abs{u-f}^2\quad\forall v\in K.
\end{equation*}
由此推出
\begin{equation*}
\abs{v-u}\le\abs{v-f}\quad\forall v\in K.
\end{equation*}
给出几何解释.
\end{exercise}

\begin{exercise}
\label{exer:5.5}
\begin{enumerate}[label=\arabic*., leftmargin=2.2em, itemsep=0.3em]
	\item 设 $(K_n)$ 为 $H$ 中不增的闭凸集序列, 使得
	      $\bigcap_nK_n\ne\emptyset$.证明对每个 $f\in H$, 序列
	      $u_n=P_{K_n}f$ (强) 收敛到一个极限, 并确定该极限.
	\item 设 $(K_n)$ 为 $H$ 中不增的非空闭凸集序列.证明对每个 $f\in H$,
	      序列 $u_n=P_{K_n}f$ (强) 收敛到一个极限, 并确定该极限.

	      设 $\varphi:H\to\R$ 为下方有界的连续函数.证明序列
	      $\alpha_n=\inf_{K_n}\varphi$ 收敛, 并确定该极限.
\end{enumerate}
\end{exercise}

\begin{exercise}[到单位球上的径向投影.]
\label{exer:5.6}
设 $E$ 为赋以范数 $\norm{\cdot}$ 的向量空间.令
\begin{equation*}
Tu=\begin{cases}
u & \text{若}\ \norm{u}\le1,\\
u/\norm{u} & \text{若}\ \norm{u}>1.
\end{cases}
\end{equation*}
\begin{enumerate}[label=\arabic*., leftmargin=2.2em, itemsep=0.3em]
	\item 证明 $\norm{Tu-Tv}\le2\norm{u-v}\quad\forall u,v\in E$.
	\item 证明一般说来常数 2 不能改进.

	      [提示: 取 $E=\R^2$, 范数 $\norm{u}=\abs{u_1}+\abs{u_2}$.]
	\item 若 $\norm{\cdot}$ 是 Hilbert 范数, 情形如何?
\end{enumerate}
\end{exercise}

\begin{exercise}[到凸锥上的投影.]
\label{exer:5.7}
设 $K\subset H$ 为顶点在原点的凸锥, 即
\begin{equation*}
0\in K\quad\text{且}\quad\lambda u+\mu v\in K
\quad\forall\lambda,\mu>0,\ \forall u,v\in K;
\end{equation*}
此外假设 $K$ 闭.

给定 $f\in H$, 证明 $u=P_Kf$ 由下述性质刻画:
\begin{equation*}
u\in K,\quad(f-u,v)\le0\quad\forall v\in K,\quad\text{且}\quad(f-u,u)=0.
\end{equation*}
\end{exercise}

\begin{exercise}
\label{exer:5.8}
设 $\Omega$ 为测度空间, $h:\Omega\to[0,+\infty)$ 为可测函数.令
\begin{equation*}
K=\{u\in L^2(\Omega);\ \abs{u(x)}\le h(x)\ \text{a.e. 在}\ \Omega\ \text{上}\}.
\end{equation*}
验证 $K$ 是 $H=L^2(\Omega)$ 中的非空闭凸集.确定 $P_K$.
\end{exercise}

\begin{exercise}
\label{exer:5.9}
设 $A\subset H$ 与 $B\subset H$ 为两个非空闭凸集, 使得 $A\cap B=\emptyset$ 且
$B$ 有界.令
\begin{equation*}
C=A-B.
\end{equation*}
\begin{enumerate}[label=\arabic*., leftmargin=2.2em, itemsep=0.3em]
	\item 证明 $C$ 闭且凸.
	\item 令 $u=P_C0$, 并把 $u$ 写成 $u=a_0-b_0$, 其中 $a_0\in A$, $b_0\in B$ (这是可能的, 因为 $u\in C$).证明
	      $\abs{a_0-b_0}=\dist(A,B)=\inf_{a\in A,b\in B}\abs{a-b}$.
	      确定 $P_Ab_0$ 与 $P_Ba_0$.
	\item 设 $a_1\in A$, $b_1\in B$ 是另一对使得
	      $\abs{a_1-b_1}=\dist(A,B)$ 的元素.证明 $u=a_1-b_1$.
	      画出一些图形, 其中对 $[a_0,b_0]$ 唯一 (相应地, 不唯一).
	\item 对 Hilbert 空间的情形, 给出 Hahn--Banach 定理第二几何形式的一个
	      简单证明.
\end{enumerate}
\end{exercise}

\begin{exercise}
\label{exer:5.10}
设 $F:H\to\R$ 为 $C^1$ 类凸函数.设 $K\subset H$ 凸, $u\in H$.证明下述性质
等价:
\begin{enumerate}[label=(\roman*), leftmargin=2.4em, itemsep=0.15em]
	\item $F(u)\le F(v)\quad\forall v\in K$,
	\item $(F'(u),v-u)\ge0\quad\forall v\in K$.
\end{enumerate}
例: $F(v)=\abs{v-f}^2$, 其中 $f\in H$ 给定.
\end{exercise}

\begin{exercise}
\label{exer:5.11}
设 $M\subset H$ 为不化为 $\{0\}$ 的闭线性子空间.设 $f\in H$, $f\notin M^\perp$.
\begin{enumerate}[label=\arabic*., leftmargin=2.2em, itemsep=0.3em]
	\item 证明
	      \begin{equation*}
	      m=\inf_{\substack{u\in M\\\abs{u}=1}}(f,u)
	      \end{equation*}
	      被唯一地达到.
	\item 设 $\varphi_1,\varphi_2,\varphi_3\in H$ 给定, $E$ 表示由
	      $\{\varphi_1,\varphi_2,\varphi_3\}$ 张成的线性空间.在下列情形确定
	      $m$:
	      \begin{enumerate}[label=(\roman*), leftmargin=2.4em, itemsep=0.1em]
	      	\item $M=E$,
	      	\item $M=E^\perp$.
	      \end{enumerate}
	\item 考察 $H=L^2(0,1)$, $\varphi_1(t)=t$, $\varphi_2(t)=t^2$, $\varphi_3(t)=t^3$ 的情形.
\end{enumerate}
\end{exercise}

\begin{exercise}[预 Hilbert 空间的完备化.]
\label{exer:5.12}
设 $E$ 为赋以标量积 $(\ ,\ )$ 的向量空间.不假设 $E$ 关于范数
$\abs{u}=(u,u)^{1/2}$ 完备 ($E$ 称为\emph{预 Hilbert 空间}).

回忆对偶空间 $E^\star$ 赋以对偶范数 $\norm{f}_{E^\star}$ 后是完备的.设
$T:E\to E^\star$ 为如下定义的映射:
\begin{equation*}
\dual{Tu}{v}_{E^\star,E}=(u,v)\quad\forall u,v\in E.
\end{equation*}
验证 $T$ 是线性等距.$T$ 是满射吗?

我们的目的是证明 $R(T)$ 在 $E^\star$ 中稠密, 且 $\norm{\cdot}_{E^\star}$ 是
Hilbert 范数.

\begin{enumerate}[label=\arabic*., leftmargin=2.2em, itemsep=0.3em]
	\item 把 $E$ 的标量积转移到 $R(T)$ 上, 并把它延拓到 $\overline{R(T)}$.
	      所得的标量积记作 $((f,g))$, 其中 $f,g\in\overline{R(T)}$.

	      验证相应的范数 $((f,f))^{1/2}$ 在 $R(T)$ 上与
	      $\norm{f}_{E^\star}$ 一致.证明
	      \begin{equation*}
	      \dual{f}{v}=((f,Tv))\quad\forall v\in E,\ \forall f\in R(T).
	      \end{equation*}
	\item 证明 $\overline{R(T)}=E^\star$.

	      [提示: 给定 $f\in E^\star$, 把 $f$ 转移为 $\overline{R(T)}$ 上的
	      线性泛函, 并使用 $\overline{R(T)}$ 中的 Riesz--Fréchet 表示定理.]

	      由此推出 $E^\star$ 关于范数 $\norm{\cdot}_{E^\star}$ 是 Hilbert
	      空间.
	\item 作出结论: $E$ 的完备化可以等同于 $E^\star$.(完备化的定义见例如
	      A.~Friedman \bref{Friedman}{3}.)
\end{enumerate}
\end{exercise}

\begin{exercise}
\label{exer:5.13}
设 $E$ 为赋以范数 $\norm{\cdot}_E$ 的向量空间.对偶范数记作
$\norm{\cdot}_{E^\star}$.回忆 (\emph{多值}) 对偶映射定义为
\begin{equation*}
F(u)=\{f\in E^\star;\ \norm{f}_{E^\star}=\norm{u}_E
\ \text{且}\ \dual{f}{u}=\norm{u}_E^2\}.
\end{equation*}
\begin{enumerate}[label=\arabic*., leftmargin=2.2em, itemsep=0.3em]
	\item 假设 $F$ 满足下述性质:
	      \begin{equation*}
	      F(u)+F(v)\subset F(u+v)\quad\forall u,v\in E.
	      \end{equation*}
	      证明范数 $\norm{\cdot}_E$ 由某个标量积导出.

	      [提示: 使用习题 5.1.]
	\item 反之, 若范数 $\norm{\cdot}_E$ 由某个标量积导出, 关于 $F$ 能说些
	      什么?

	      [提示: 使用习题 5.12 与 1.1.]
\end{enumerate}
\end{exercise}

\begin{exercise}
\label{exer:5.14}
设 $a:H\times H\to\R$ 为双线性连续形式, 使得
\begin{equation*}
a(v,v)\ge0\quad\forall v\in H.
\end{equation*}
证明函数 $v\mapsto F(v)=a(v,v)$ 凸, 属于 $C^1$ 类, 并确定它的微分.
\end{exercise}

\begin{exercise}
\label{exer:5.15}
设 $G\subset H$ 为 Hilbert 空间 $H$ 的线性子空间; $G$ 赋以 $H$ 的范数.设
$F$ 为 Banach 空间.设 $S:G\to F$ 为有界线性算子.证明存在有界线性算子
$T:H\to F$, 它延拓 $S$, 且使得
\begin{equation*}
\norm{T}_{\mathcal{L}(H,F)}=\norm{S}_{\mathcal{L}(G,F)}.
\end{equation*}
\end{exercise}

\begin{exercise}[三元组 $V\subset H\subset V^\star$.]
\label{exer:5.16}
设 $H$ 为赋以标量积 $(\ ,\ )$ 及相应范数 $\abs{\cdot}$ 的 Hilbert 空间.设
$V\subset H$ 为在 $V$ 中稠密的线性子空间.
% 原书此处印作 "dense in V", 按上下文 (并与 5.2 节注记 3 的同一表述对照)
% 应为 "dense in H". 依技能规范保留原样并记录, 不擅自更正.
假设 $V$ 有自己的范数
$\norm{\cdot}$, 且 $V$ 关于 $\norm{\cdot}$ 是 Banach 空间.还假设嵌入
$V\subset H$ 连续, 即 $\abs{v}\le C\norm{v}\ \forall v\in V$.考虑由
\begin{equation*}
\dual{Tu}{v}_{V^\star,V}=(u,v)\quad\forall u\in H,\ \forall v\in V
\end{equation*}
定义的算子 $T:H\to V^\star$.
\begin{enumerate}[label=\arabic*., leftmargin=2.2em, itemsep=0.3em]
	\item 证明 $\norm{Tu}_{V^\star}\le C\abs{u}\ \forall u\in H$.
	\item 证明 $T$ 是单射.
	\item 证明若 $V$ 自反, 则 $R(T)$ 在 $V^\star$ 中稠密.
	\item 给定 $f\in V^\star$, 证明 $f\in R(T)$ 当且仅当存在常数 $a\ge0$,
	      使得 $\abs{\dual{f}{v}_{V^\star,V}}\le a\abs{v}\ \forall v\in V$.
\end{enumerate}
\end{exercise}

\begin{exercise}
\label{exer:5.17}
设 $M,N\subset H$ 为两个闭线性子空间.假设
$(u,v)=0\ \forall u\in M,\ \forall v\in N$.证明 $M+N$ 闭.
\end{exercise}

\begin{exercise}
\label{exer:5.18}
设 $E$ 为 Banach 空间, $H$ 为 Hilbert 空间.设
$T\in\mathcal{L}(E,H)$.证明下述性质等价:
\begin{enumerate}[label=(\roman*), leftmargin=2.4em, itemsep=0.15em]
	\item $T$ 有左逆,
	\item 存在常数 $C$, 使得 $\norm{u}\le C\abs{Tu}\ \forall u\in E$.
\end{enumerate}
\end{exercise}

\begin{exercise}
\label{exer:5.19}
设 $(u_n)$ 为 $H$ 中的序列, 使得 $u_n\weakto u$ 弱收敛.假设
$\limsup\abs{u_n}\le\abs{u}$.不依赖命题  \ref{prop:3.32} 证明
$u_n\to u$ 强收敛.
\end{exercise}

\begin{exercise}
\label{exer:5.20}
假设 $S\in\mathcal{L}(H)$ 满足 $(Su,u)\ge0\ \forall u\in H$.
\begin{enumerate}[label=\arabic*., leftmargin=2.2em, itemsep=0.3em]
	\item 证明 $N(S)=R(S)^\perp$.
	\item 证明对每个 $t>0$, $I+tS$ 是双射.
	\item 证明
	      \begin{equation*}
	      \lim_{t\to+\infty}(I+tS)^{-1}f=P_{N(S)}f\quad\forall f\in H.
	      \end{equation*}
	      [提示: 有两种方法:
	      \begin{enumerate}[label=(\alph*), leftmargin=2.2em, itemsep=0.1em]
	      	\item 分别考虑 $f\in N(S)$ 与 $f\in R(S)$ 的情形.
	      	\item 使用弱收敛.
	      \end{enumerate}]
\end{enumerate}
\end{exercise}

\begin{exercise}[线性压缩的迭代. Kakutani--Yosida 遍历定理.]
\label{exer:5.21}
设 $T\in\mathcal{L}(H)$ 满足 $\norm{T}\le1$.给定 $f\in H$ 与整数 $n\ge1$,
令
\begin{equation*}
\sigma_n(f)=\frac1n\bigl(f+Tf+T^2f+\cdots+T^{n-1}f\bigr)
\end{equation*}
以及
\begin{equation*}
\mu_n(f)=\Bigl(\frac{I+T}{2}\Bigr)^nf.
\end{equation*}
我们的目的是证明
\begin{equation*}
\lim_{n\to\infty}\sigma_n(f)=\lim_{n\to\infty}\mu_n(f)=P_{N(I-T)}f.
\end{equation*}
\begin{enumerate}[label=\arabic*., leftmargin=2.2em, itemsep=0.3em]
	\item 验证 $N(I-T)=R(I-T)^\perp$.
	\item 假设 $f\in R(I-T)$.证明存在常数 $C$, 使得
	      $\abs{\sigma_n(f)}\le C/n\ \forall n\ge1$.
	\item 由此推出: 对每个 $f\in H$, 有
	      \begin{equation*}
	      \lim_{n\to\infty}\sigma_n(f)=P_{N(I-T)}f.
	      \end{equation*}
	\item 令 $S=\frac12(I+T)$.证明
	      \begin{equation}
	      \abs{u-Su}^2+\abs{Su}^2\le\abs{u}^2\quad\forall u\in H.
	      \tag{1}
	      \end{equation}
	      由此推出
	      \begin{equation*}
	      \sum_{i=0}^\infty\abs{S^iu-S^{i+1}u}^2\le\abs{u}^2\quad\forall u\in H
	      \end{equation*}
	      以及
	      \begin{equation*}
	      \abs{S^n(u-Su)}\le\frac{\abs{u}}{\sqrt{n+1}}
	      \quad\forall u\in H,\ \forall n\ge1.
	      \end{equation*}
	\item 假设 $f\in R(I-T)$.证明存在常数 $C$, 使得
	      $\abs{\mu_n(f)}\le C/\sqrt{n}\ \forall n\ge1$.
	\item 由此推出: 对每个 $f\in H$, 有
	      \begin{equation*}
	      \lim_{n\to\infty}\mu_n(f)=P_{N(I-T)}f.
	      \end{equation*}
\end{enumerate}
\end{exercise}

\begin{exercise}
\label{exer:5.22}
设 $C\subset H$ 为非空闭凸集, $T:C\to C$ 为非线性压缩, 即
\begin{equation*}
\abs{Tu-Tv}\le\abs{u-v}\quad\forall u,v\in C.
\end{equation*}
\begin{enumerate}[label=\arabic*., leftmargin=2.2em, itemsep=0.3em]
	\item 设 $(u_n)$ 为 $C$ 中的序列, 使得
	      \begin{equation*}
	      u_n\weakto u\ \text{弱收敛}\quad\text{且}\quad(u_n-Tu_n)\to f
	      \ \text{强收敛}.
	      \end{equation*}
	      证明 $u-Tu=f$.

	      [提示: 先考虑 $C=H$ 的情形, 并使用不等式
	      $((u-Tu)-(v-Tv),u-v)\ge0\ \forall u,v$.]
	\item 由此推出: 若 $C$ 有界且 $T(C)\subset C$, 则 $T$ 有不动点.

	      [提示: 考虑 $T_\varepsilon u=(1-\varepsilon)Tu+\varepsilon a$,
	      其中 $a\in C$ 固定, $\varepsilon>0$, $\varepsilon\to0$.]
\end{enumerate}
\end{exercise}

\begin{exercise}[Zarantonello 不等式.]
\label{exer:5.23}
设 $T:H\to H$ 为 (非线性) 压缩.假设 $\alpha_1,\alpha_2,\dots,\alpha_n\in\R$
满足 $\alpha_i\ge0\ \forall i$ 且 $\sum_{i=1}^n\alpha_i=1$.假设
$u_1,u_2,\dots,u_n\in H$, 并令
\begin{equation*}
\sigma=\sum_{i=1}^n\alpha_iu_i.
\end{equation*}
证明
\begin{equation*}
\abs{T\sigma-\sum_{i=1}^n\alpha_iTu_i}^2
\le\frac12\sum_{i,j=1}^n\alpha_i\alpha_j
\bigl(\abs{u_i-u_j}^2-\abs{Tu_i-Tu_j}^2\bigr).
\end{equation*}
[提示: 写出
\begin{equation*}
\abs{T\sigma-\sum_{i=1}^n\alpha_iTu_i}^2
=\sum_{i,j=1}^n\alpha_i\alpha_j(T\sigma-Tu_i,T\sigma-Tu_j),
\end{equation*}
并使用恒等式 $(a,b)=\frac12\bigl(\abs{a}^2+\abs{b}^2-\abs{a-b}^2\bigr)$.]

当 $T$ 是等距 (即 $\abs{Tu-Tv}=\abs{u-v}\ \forall u,v\in H$) 时能推出什么?
\end{exercise}

\begin{exercise}[Banach--Saks 性质.]
\label{exer:5.24}
\begin{enumerate}[label=\arabic*., leftmargin=2.2em, itemsep=0.3em]
	\item 假设 $(u_n)$ 为 $H$ 中的序列, 使得 $u_n\weakto0$ 弱收敛.用归纳法
	      构造子列 $(u_{n_j})$, 使得 $u_{n_1}=u_1$ 且
	      \begin{equation*}
	      \abs{(u_{n_j},u_{n_k})}\le\frac1k
	      \quad\forall k\ge2,\ \forall j=1,2,\dots,k-1.
	      \end{equation*}
	      由此推出: 由 $\sigma_p=\frac1p\sum_{j=1}^pu_{n_j}$ 定义的序列
	      $(\sigma_p)$ 当 $p\to\infty$ 时强收敛于 0.

	      [提示: 估计 $\abs{\sigma_p}^2$.]
	\item 假设 $(u_n)$ 为 $H$ 中的有界序列.证明存在子列 $(u_{n_j})$, 使得
	      序列 $\sigma_p=\frac1p\sum_{j=1}^pu_{n_j}$ 当 $p\to\infty$ 时强
	      收敛到一个极限.

	      与推论  \ref{cor:3.8} 及习题 3.4 比较.
\end{enumerate}
\end{exercise}

\begin{exercise}[Opial 引理的若干变体.]
\label{exer:5.25}
设 $K\subset H$ 为非空闭凸集.设 $(u_n)$ 为 $H$ 中的序列, 使得对每个
$v\in K$, 序列 $(\abs{u_n-v})$ 不增.
\begin{enumerate}[label=\arabic*., leftmargin=2.2em, itemsep=0.3em]
	\item 验证序列 $(\dist(u_n,K))$ 不增.
	\item 证明序列 $(P_Ku_n)$ 强收敛到一个极限, 记作 $\ell$.

	      [提示: 使用习题 5.4.]
	\item 这里假设序列 $(u_n)$ 具有性质
	      \begin{equation*}
	      \text{(P)}\qquad
	      \begin{cases}
	      \text{每当某个子列 }(u_{n_k})\text{ 弱收敛到某个极限 }u\in H\text{ 时,}\\
	      \text{都有 }u\in K.
	      \end{cases}
	      \end{equation*}
	      证明 $u_n\weakto\ell$ 弱收敛.
	\item 这里假设 $\bigcup_{\lambda>0}\lambda(K-K)=H$.证明存在某个 $u\in H$,
	      使得 $u_n\weakto u$ 弱收敛且 $P_Ku=\ell$.
	\item 这里假设 $\Int K\ne\emptyset$.证明存在某个 $u\in H$, 使得
	      $u_n\to u$ 强收敛.

	      [提示: 先考虑 $K$ 是单位球的情形, 再考虑一般情形.]
	\item 令 $\sigma_n=\frac1n(u_1+u_2+\cdots+u_n)$, 并假设序列 $(\sigma_n)$
	      满足性质 (P).证明 $\sigma_n\weakto\ell$ 弱收敛.
\end{enumerate}
\end{exercise}

\begin{exercise}
\label{exer:5.26}
假设 $(e_n)$ 是 $H$ 的正交基.
\begin{enumerate}[label=\arabic*., leftmargin=2.2em, itemsep=0.3em]
	\item 验证 $e_n\weakto0$ 弱收敛.

	      设 $(a_n)$ 为 $\R$ 中的有界序列, 令
	      $u_n=\frac1n\sum_{i=1}^na_ie_i$.
	\item 证明 $\abs{u_n}\to0$.
	\item 证明 $\sqrt{n}\,u_n\weakto0$ 弱收敛.
\end{enumerate}
\end{exercise}

\begin{exercise}
\label{exer:5.27}
设 $D\subset H$ 为子集, 使得由 $D$ 张成的线性空间在 $H$ 中稠密.设
$(E_n)_{n\ge1}$ 为 $H$ 中两两正交的闭子空间序列.假设
\begin{equation*}
\sum_{n=1}^\infty\abs{P_{E_n}u}^2=\abs{u}^2\quad\forall u\in D.
\end{equation*}
证明 $H$ 是诸 $E_n$ 的 Hilbert 和.
\end{exercise}

\begin{exercise}
\label{exer:5.28}
假设 $H$ 可分.
\begin{enumerate}[label=\arabic*., leftmargin=2.2em, itemsep=0.3em]
	\item 设 $V\subset H$ 为在 $H$ 中稠密的线性子空间.证明 $V$ 包含 $H$ 的
	      一个正交基.
	\item 设 $(e_n)_{n\ge1}$ 为 $H$ 中的正交序列, 即 $(e_i,e_j)=\delta_{ij}$.
	      证明存在 $H$ 的一个正交基, 它包含 $\bigcup_{n=1}^\infty\{e_n\}$.
\end{enumerate}
\end{exercise}

\begin{exercise}[Grothendieck 的一个引理.]
\label{exer:5.29}
设 $\Omega$ 为 $\abs{\Omega}<\infty$ 的测度空间.设 $E$ 为 $L^p(\Omega)$
($1\le p<\infty$) 的闭子空间.假设 $E\subset L^\infty(\Omega)$.我们的目的
是证明 $\dim E<\infty$.
\begin{enumerate}[label=\arabic*., leftmargin=2.2em, itemsep=0.3em]
	\item 证明存在常数 $C$, 使得
	      \begin{equation*}
	      \norm{u}_\infty\le C\norm{u}_p\quad\forall u\in E.
	      \end{equation*}
	      [提示: 使用推论~
\ref{cor:2.8}.]
	\item 证明存在常数 $M$, 使得
	      \begin{equation*}
	      \norm{u}_\infty\le M\norm{u}_2\quad\forall u\in E.
	      \end{equation*}
	      [提示: 区分 $1\le p\le2$ 与 $2<p<\infty$ 两种情形.]
	\item 由此推出 $E$ 是 $L^2(\Omega)$ 的闭子空间.

	      下文假设 $\dim E=\infty$.设 $(e_n)_{n\ge1}$ 为 $E$ (赋以 $L^2$
	      标量积) 中的正交序列.
	\item 固定任一整数 $k\ge1$.证明存在零测集 $\omega\subset\Omega$, 使得
	      \begin{equation*}
	      \Bigl|\sum_{i=1}^k\alpha_ie_i(x)\Bigr|
	      \le M\Bigl(\sum_{i=1}^k\alpha_i^2\Bigr)^{1/2}
	      \quad\forall x\in\Omega\setminus\omega,\
	      \forall\alpha=(\alpha_1,\alpha_2,\dots,\alpha_k)\in\R^k.
	      \end{equation*}
	      [提示: 先考虑 $\alpha\in\Q^k$ 的情形.]
	\item 由此推出 $\sum_{i=1}^k\abs{e_i(x)}^2\le M^2
	      \ \forall x\in\Omega\setminus\omega$.
	\item 作出结论.
\end{enumerate}
\end{exercise}

\begin{exercise}
\label{exer:5.30}
设 $(e_n)_{n\ge1}$ 为 $H=L^2(0,1)$ 中的正交序列.设 $p(t)$ 为 $H$ 中给定
的函数.
\begin{enumerate}[label=\arabic*., leftmargin=2.2em, itemsep=0.3em]
	\item 证明对每个 $t\in[0,1]$, 有
	      \begin{equation}
	      \sum_{n=1}^\infty\Bigl|\int_0^tp(s)e_n(s)\,ds\Bigr|^2
	      \le\int_0^t\abs{p(s)}^2\,ds.
	      \tag{1}
	      \end{equation}
	\item 由此推出
	      \begin{equation}
	      \sum_{n=1}^\infty\int_0^1\Bigl|\int_0^tp(s)e_n(s)\,ds\Bigr|^2dt
	      \le\int_0^1\abs{p(t)}^2(1-t)\,dt.
	      \tag{2}
	      \end{equation}
	\item 现在假设 $(e_n)_{n\ge1}$ 是 $H$ 的正交基.证明 (1) 与 (2) 变成
	      等式.
	\item 反之, 假设 (2) 中等号成立, 且 $p(t)\ne0$ 几乎处处.证明
	      $(e_n)_{n\ge1}$ 是正交基.
\end{enumerate}
例: $p\equiv1$.
\end{exercise}

\begin{exercise}[Haar 基.]
\label{exer:5.31}
给定整数 $n\ge1$, 把它写成 $n=k+2^p$, 其中 $p\ge0$, $k\ge0$ 是由条件
$k\le2^p-1$ 唯一确定的整数.考虑由下式定义在 $(0,1)$ 上的函数
\begin{equation*}
\varphi_n(t)=\begin{cases}
2^{p/2} & \text{若}\ k2^{-p}<t<\bigl(k+\frac12\bigr)2^{-p},\\
-2^{p/2} & \text{若}\ \bigl(k+\frac12\bigr)2^{-p}<t<(k+1)2^{-p},\\
0 & \text{其他情形}.
\end{cases}
\end{equation*}
令 $\varphi_0\equiv1$, 证明 $(\varphi_n)_{n\ge0}$ 是 $L^2(0,1)$ 的正交基.
\end{exercise}

\begin{exercise}[Rademacher 系与 Walsh 基.]
\label{exer:5.32}
对每个整数 $i\ge0$, 考虑由 $r_i(t)=(-1)^{[2^it]}$ 定义在 $(0,1)$ 上的
函数 $r_i(t)$ (如通常那样, $[x]$ 表示不超过 $x$ 的最大整数).
\begin{enumerate}[label=\arabic*., leftmargin=2.2em, itemsep=0.3em]
	\item 验证 $(r_i)_{i\ge0}$ 是 $L^2(0,1)$ 中的正交序列 (称为
	      \emph{Rademacher 系}).
	\item $(r_i)_{i\ge0}$ 是正交基吗?

	      [提示: 考虑函数 $u=r_1r_2$.]
	\item 给定整数 $n\ge0$, 考虑它的二进制表示
	      $n=\sum_{i=0}^\ell\alpha_i2^i$, 其中 $\alpha_i\in\{0,1\}$.令
	      \begin{equation*}
	      w_n(t)=\prod_{i=0}^\ell r_{i+1}(t)^{\alpha_i}.
	      \end{equation*}
	      证明 $(w_n)_{n\ge0}$ 是 $L^2(0,1)$ 的正交基 (称为
	      \emph{Walsh 基}).注意 $(r_i)_{i\ge0}$ 是 $(w_n)_{n\ge0}$ 的
	      子集.
\end{enumerate}
\end{exercise}
